\section{结论}\label{sec:conclusion}

本文综述了 ABC 猜想从提出至今的完整脉络。概括而言：

\begin{enumerate}[label=(\arabic*)]
    \item \textbf{历史主线}是“类比先行、翻译跟进”：Mason--Stothers（1981）在函数域中证明了指数为零的精确定理，Masser--Oesterl\'e（1985）经曲线--数域平行论把它翻译为整数世界的猜想，Frey 曲线把猜想与 Szpiro、费马、Mordell 织成等价网络，Granville--Tucker（2002）将其纲领化为数论的“新常态”。
    \item \textbf{方法层面}，三条战线各自清晰：无条件下界沿 Baker 对数线性型理论推进到指数 $1/3$（Stewart--Tijdeman、Stewart--Yu），与目标的指数 $1$ 之间存在方法论鸿沟；反方向上，Elkies 的 Belyi 构造与 \textit{abc@home} 的搜索持续逼近猜想的尖锐度边界；进攻方向上，望月新的 IUT（2012）宣称穿越这条鸿沟，但其核心推论 3.12 自 Scholze--Stix（2018）的质疑以来处于未决争议，ABC 在数学共同体的通行标准下仍是开放猜想。
    \item \textbf{后果层面}，ABC 是性价比最高的“总开关”：渐近费马、广义费马的有效有限性、Catalan 型方程的分类、Hall 弱形式、平方自由值、有效 Mordell 与 Szpiro 控制——几乎所有指数型丢番图问题的“$\varepsilon$ 版本”都是它的直接推论。
    \item \textbf{前沿方向}包括 IUT 缺口的解决或绕越、显式 $K_\varepsilon$ 与质量分布的定量理论、特征 $p$ 类比的正确表述、Vojta 型高维推广，以及大型证明可传递性的制度化与形式化。
\end{enumerate}

作为“加法与乘法互不相容”这一朴素信念的精确化，ABC 猜想浓缩了数论的根本张力：乘法世界由素数各向同性地支配，加法世界却要求结构选择，而两者的对话被一个 $\varepsilon$ 隔开。一个世纪的函数域定理、二十年的显式计算与一场至今未决的证明争议，都在测量这个 $\varepsilon$ 的宽度。无论它最终由何条路线闭合——IUT 的修复、Arakelov 几何的新不等式，或尚未出现的语言——其闭合之日，都将是丢番图方程版图的重新绘制之时。
